% CV académico (español).
% Compilar desde la raíz del repositorio con:  make carolina_branas_cv_espanol

\documentclass{resume}
\usepackage[left=0.6in,top=0.5in,right=0.6in,bottom=0.45in]{geometry}
\usepackage{vwcol}
\usepackage{natbib}
\setlength{\bibsep}{2.5pt}
\usepackage[svgnames]{xcolor}
\usepackage{hyperref}
\usepackage{graphicx}
\hypersetup{%
  colorlinks=true,
  allcolors=DarkRed,
  pdfborderstyle={/S/U/W 1}
}

% espaciado de secciones compacto
\titlespacing{\section}{0pt}{6pt}{3pt}

\begin{document}

%%%%%%%%%%%%%%%%%%%%%%%%%%%%%%%%%%%%%%%%%%%%%%%%% TÍTULO
\vspace*{-2em}

\begin{minipage}[t]{0.76\textwidth}
    \raggedright
    {\Huge \textbf{Carolina Brañas}}\\[0.3em]
    {\normalsize \textbf{Científica de Datos}}\\[0.5em]
    {\footnotesize
    A Coruña, España $\bullet$
    (+34) 644 004 477 $\bullet$
    \href{mailto:carobrasor@gmail.com}{carobrasor@gmail.com} $\bullet$
    \href{https://carobs9.github.io/}{carobs9.github.io} \\[0.2em]
    \href{https://github.com/carobs9}{github.com/carobs9} $\bullet$
    \href{https://www.linkedin.com/in/carolinabranas/}{linkedin.com/in/carolinabranas} $\bullet$
    \href{https://huggingface.co/carobs9}{huggingface.co/carobs9}
    }\\[0.7em]
    {\small
    Científica de datos con experiencia práctica en modelos de aprendizaje automático y
    procesamiento del lenguaje natural (NLP), diseño de flujos de datos escalables y
    conversión de datos complejos en información accionable.
    Intereses: aprendizaje automático, análisis de redes, NLP y datos geoespaciales.
    }
\end{minipage}%
\hfill
\begin{minipage}[t]{0.22\textwidth}
    \vspace*{0.2em}
    \begin{flushright}
        \includegraphics[width=2.7cm,clip,trim=0 0 0 0]{profile.jpg}
    \end{flushright}
\end{minipage}
\vspace{0.2em}

%%%%%%%%%%%%%%%%%%%%%%%%%%%%%%%%%%%%%%%%%%%%%%%%% EXPERIENCIA
\section{EXPERIENCIA}
\begin{content}

    \begin{position}{Asistente de Investigación, Data Manager}{Feb 2026 -- Presente}{Universidad de Copenhague, Dinamarca}{Prof.~Asger Mose Wingender}{}
        \item Contribuyo al proyecto {\href{https://www.carlsbergfondet.dk/en/what-we-have-funded/cf25-1677/}{Cataloguing Crop Traits and Breeders Across the OECD 1980-2025}}, financiado por la Carlsberg Foundation.
        \item Gestiono la creación del primer repositorio unificado de variedades de cereal registradas en la OCDE desde 1980, con sus características, historial de mejora genética y las empresas que las desarrollaron.
        \item Gestiono a un estudiante de apoyo en el proyecto, coordinando y revisando su trabajo.
    \end{position}

    \begin{position}{Asistente de Investigación, Ciencia de Datos}{Mayo 2024 -- Dic 2025}{Universidad de Copenhague, Dinamarca}{Prof.~Jeanet Bentzen}{}
        \item Contribuí al proyecto {\href{https://www.economics.ku.dk/research/externally-funded-research_new/shocking-religion/}{Shocking Religion}} sobre el impacto económico de la religión, dentro de un equipo multidisciplinar.
        \item Desarrollé modelos de temas para descubrir tendencias temáticas en datos textuales y flujos basados en LLM y RAG para la extracción de información en documentos.
        \item Calculé distancias desde distintas ubicaciones hasta posibles zonas de peligro sísmico para cada uno de los cuatro niveles de riesgo (1--4), con Rasterio, GeoPandas y QGIS.
        \item Diseñé y contenedoricé flujos de datos escalables (Docker); gestioné la ingesta y el preprocesamiento a gran escala; desplegué en infraestructura en la nube (UCloud).
    \end{position}

    \begin{position}{Científica de Datos}{Oct 2023 -- Mayo 2024}{Above Sports, Dinamarca}{}{}
        \item Desarrollé modelos de computer vision para la detección de logotipos de marca en vídeo.
        \item Automaticé y contenedoricé flujos de datos, y colaboré con los equipos de producto para refinar la calidad de los resultados.
    \end{position}

    \begin{position}{Estratega de Marketing}{Sep 2021 -- Mayo 2022}{Crescendo Collective, Estados Unidos}{}{}
        \item Analicé datos de campañas (Google Analytics, Google Ads) y automaticé informes internos con Python.
        \item Realicé análisis de competencia y de audiencias con el equipo de datos; presenté informes a clientes.
    \end{position}

\end{content}

%%%%%%%%%%%%%%%%%%%%%%%%%%%%%%%%%%%%%%%%%%%%%%%%% EDUCACIÓN
\section{EDUCACIÓN}
\begin{content}
    {\bf M.Sc.\ en Ciencia de Datos Sociales} \hfill {\bf 2022 -- 2024} \\
    Universidad de Copenhague, Dinamarca \\
    {\bf \em Tesis:} {\href{https://github.com/carobs9/segregation-madrid}{Movilidad y segregación por ingresos en Madrid, España.}} \\
    {\small {\em Optativas:} Aprendizaje Automático Avanzado, Ciencia de Datos Geoespaciales, Ciencia de Redes Avanzada (ITU); Procesamiento del Lenguaje Natural (DIKU).} \\[0.4em]
    {\bf B.Sc.\ en Marketing, Minor en Ingeniería Informática} \hfill {\bf 2018 -- 2022} \\
    Cardinal Stritch University, Estados Unidos \\
    {\small {\em Honores:} Magna Cum Laude, Lista del Decano (2018--2022), mejor promedio del grado en Marketing (2022).}
\end{content}

%%%%%%%%%%%%%%%%%%%%%%%%%%%%%%%%%%%%%%%%%%%%%%%%% HABILIDADES
\section{HABILIDADES}
\begin{content}
    {\bf Programación y Datos:} Python {\footnotesize (pandas, numpy, pytorch, tensorflow, scikit-learn)}, SQL, Bash, \LaTeX \\
    {\bf Aprendizaje Automático y NLP:} Transformers, topic modeling {\footnotesize (UMAP, HDBSCAN, BERTopic)}, Hugging Face, Weights \& Biases, deep learning, modelos predictivos, feature engineering, computer vision \\
    {\bf Análisis, Visualización y Geoespacial:} análisis estadístico, tratamiento de datos, Matplotlib, Seaborn, Tableau, Power BI, Geopandas, Rasterio, QGIS \\
    {\bf Cloud y DevOps:} Docker, Git, Linux, AWS, Azure, UCloud, CI/CD \\
    {\bf Idiomas:} español {\footnotesize (nativo)}, gallego {\footnotesize (nativo)}, inglés {\footnotesize (competencia profesional)}, danés {\footnotesize (principiante)}
\end{content}

%%%%%%%%%%%%%%%%%%%%%%%%%%%%%%%%%%%%%%%%%%%%%%%%% PREMIOS
\section{PREMIOS}
\begin{content}
    \prize[s]{Lista del Decano}{2018--2022}{{\small todos los semestres del grado, por un promedio superior a 3.5/4.0}} \\
    \prize[s]{Mejor Promedio del B.Sc.\ en Marketing}{2022}{{\small promedio más alto de la promoción 2022}} \\
    \prize[s]{Beca Académica y Deportiva}{2018--2022}{{\small beca completa de cuatro años por excelencia académica y fútbol}}
\end{content}

\end{document}
